\documentclass[a4paper,12pt]{article} 

\usepackage[spanish]{babel}
\usepackage[utf8]{inputenc}
\usepackage[T1]{fontenc}

\usepackage{geometry}
\geometry{margin=2.5cm}

\usepackage{hyperref}
\usepackage{microtype}
\usepackage{color, soul}

\usepackage{fancyhdr}

\pagestyle{fancy}
\fancyhf{}
\lhead{Guía de Ejercicios 1.1}
\rhead{Tecnología Digital VII}
\cfoot{\thepage}

\title{\vspace{-2cm}\textbf{Guía de Ejercicios 1.1}\\[0.2cm]
\large Tecnología Digital VII}
\author{Juan Ignacio Elosegui}
\date{Universidad Torcuato Di Tella \\ Marzo 2026}

\begin{document}
    \maketitle
    \section*{Ejercicio 1}
        \subsection*{Inciso 1}
            \textbf{Entidad: Cliente}
                \begin{itemize}
                    \item \underline{\texttt{id\_cliente}} (PK)
                \end{itemize}

            \textbf{Relación: tiene}
                \begin{itemize}
                    \item Cliente (1) --- (1) Cuenta
                \end{itemize}

            \textbf{Entidad: Cuenta bancaria}
                \begin{itemize}
                    \item \underline{\texttt{id\_cuenta}} (PK)
                \end{itemize}

       \subsection*{Inciso 2}
            \textbf{Entidad: Cliente}
                \begin{itemize}
                    \item \underline{\texttt{id\_cliente}} (PK)
                \end{itemize}

            \textbf{Relación: tiene}
                \begin{itemize}
                    \item Cliente (1) --- (N)
                \end{itemize}

            \textbf{Entidad: Cuenta bancaria}
                \begin{itemize}
                    \item \underline{\texttt{id\_cuenta}} (PK)
                \end{itemize}
            
        \subsection*{Inciso 3}
            \textbf{Entidad: Cliente}
                \begin{itemize}
                    \item \underline{\texttt{id\_cliente}} (PK)
                \end{itemize}

            \textbf{Relación: tiene}
                \begin{itemize}
                    \item Cliente (1) --- Cuentas (N)
                \end{itemize}

            \textbf{Entidad: Cuenta bancaria}
                \begin{itemize}
                    \item \underline{\texttt{id\_cuenta}} (PK)
                \end{itemize}
                \noindent El valor de \texttt{id\_cuenta} puede ser \texttt{NULL}.

        \subsection*{Inciso 4}
            \textbf{Entidad: Cuenta}
                \begin{itemize}
                    \item \underline{\texttt{id\_cuenta}} (PK)
                \end{itemize}

            \textbf{Relación: asociada a}
                \begin{itemize}
                    \item Cuentas (N) --- Clientes (N)
                \end{itemize}

            \textbf{Entidad: Clientes}
                \begin{itemize}
                    \item \texttt{id\_de\_dueños}
                \end{itemize}
                \noindent \hl{Importa decir ``cuentas bancarias a más de un cliente'' o ``clientes a más de una cuenta bancaria''? No sé si está bien la lógica.}

        \subsection*{Inciso 5}
            \textbf{Entidad: Postulante}
                \begin{itemize}
                    \item \underline{\texttt{CUIT}} (PK)
                    \item \texttt{apellido}
                    \item \texttt{nombre}
                    \item \texttt{domicilio}
                    \item \texttt{telefonos}
                    \item \texttt{direcciones\_de\_mail}
                \end{itemize}

            \textbf{Relación: tiene}
                \begin{itemize}
                    \item Postulante (1) --- (N) Habilidades
                \end{itemize}

            \textbf{Entidad: Habilidades}
                \begin{itemize}
                    \item \underline{\texttt{CUIT}} (PK)
                    \item \texttt{habilidades}
                \end{itemize}
                \noindent \hl{¿La clave primaria de las dos entidades deben ser iguales? Si no, ¿cómo sabe que estamos hablando del mismo postulante?}

        \subsection*{Inciso 6}
            \textbf{Entidad: }
                \begin{itemize}
                    \item \underline{\texttt{}} (PK)
                    \item ...
                \end{itemize}

            \textbf{Relación: }
                \begin{itemize}
                    \item 
                \end{itemize}

            \textbf{Entidad: }
                \begin{itemize}
                    \item \underline{\texttt{}} (PK)
                    \item ...
                \end{itemize}

        \subsection*{Inciso 7}
            \textbf{Entidad: }
                \begin{itemize}
                    \item \underline{\texttt{}} (PK)
                    \item ...
                \end{itemize}

            \textbf{Relación: }
                \begin{itemize}
                    \item 
                \end{itemize}

            \textbf{Entidad: }
                \begin{itemize}
                    \item \underline{\texttt{}} (PK)
                    \item ...
                \end{itemize}

        \subsection*{Inciso 8}
            \textbf{Entidad: }
                \begin{itemize}
                    \item \underline{\texttt{}} (PK)
                    \item ...
                \end{itemize}

            \textbf{Relación: }
                \begin{itemize}
                    \item 
                \end{itemize}

            \textbf{Entidad: }
                \begin{itemize}
                    \item \underline{\texttt{}} (PK)
                    \item ...
                \end{itemize}

    \section*{Ejercicio 2}

    \section*{Ejercicio 3}
        \subsection*{Inciso (a)}
        \subsection*{Inciso (b)}
        \subsection*{Inciso (c)}
        \subsection*{Inciso (d)}
        \subsection*{Inciso (e)}
        \subsection*{Inciso (f)}
        \subsection*{Inciso (g)}
        \subsection*{Inciso (h)}
        \subsection*{Inciso (i)}
        \subsection*{Inciso (j)}
        \subsection*{Inciso (k)}
        \subsection*{Inciso (l)}

    \section*{Ejercicio 4}
        \subsection*{Inciso 1}
            \subsubsection*{Inciso (a)}
            \subsubsection*{Inciso (b)}
            \subsubsection*{Inciso (c)}
        \subsection*{Inciso 2}
            \subsubsection*{Inciso (a)}
            \subsubsection*{Inciso (b)}
            \subsubsection*{Inciso (c)}

    \section*{Ejercicio 5}
        \subsection*{Inciso 1}
            \subsubsection*{Inciso (a)}
            \subsubsection*{Inciso (b)}
            \subsubsection*{Inciso (c)}
            \subsubsection*{Inciso (d)}
            \subsubsection*{Inciso (e)}
            \subsubsection*{Inciso (f)}
            \subsubsection*{Inciso (g)}
            \subsubsection*{Inciso (h)}
            \subsubsection*{Inciso (i)}
            \subsubsection*{Inciso (j)}
        \subsection*{Inciso 2}

    \section*{Ejercicio 6}
        \subsection*{Inciso (a)}
        \subsection*{Inciso (b)}
        \subsection*{Inciso (c)}
        \subsection*{Inciso (d)}

    \section*{Ejercicio 7}
        \subsection*{Inciso (a)}
        \subsection*{Inciso (b)}
        \subsection*{Inciso (c)}
        
\end{document} 
